% Generated by code/08_generate_results_tables.R
\begin{table}[H]
\centering
\scriptsize
\caption{Stage 3 behavioural identification design}
\label{tab:app_stage3_behavioural_design}
\resizebox{\textwidth}{!}{%
\begin{tabular}{@{}L{2.7cm}L{4.1cm}L{3.7cm}L{5.1cm}@{}}
\toprule
Component & Baseline treatment & Identification status & Data for structural extension \\
\midrule
Reallocation share & Estimated from the CbCR fixed-effect semi-elasticity and applied only where a legal collection channel remains active. & Identified as a reduced-form profit response, not as a structural destination-choice parameter. & Private firm-level CbCR or unconsolidated affiliate accounts observed around tax shocks, with group identifiers and affiliate balance sheets. \\
Low etr nonlinearity & Quadratic and logarithmic semi-elasticity specifications are estimated, but the preferred estimator is still chosen by BIC. A lower-tail support diagnostic reports below-5 per cent observations and an illustrative tail-balanced reweighting design. & The current public CbCR panel identifies tax responsiveness more credibly than it identifies the curvature of the response at the extreme lower ETR tail. Tail upweighting is not treated as IPW because the public data do not observe a sampling or treatment-selection probability for below-5 per cent ETR cells. & MNE-jurisdiction microdata or a quasi-experimental tax shock with dense support below 5 per cent, plus a defensible model for the probability that an observation enters the lower tail. \\
Reallocation cap & No discretionary cap is imposed; the response is bounded only to the unit interval. & A cap below one is not identified from public aggregate CbCR data. & Observed post-reform relocation shares from private CbCR or tax-return microdata. \\
Sesh notch bunching & Implemented as an upper-bound liability filter inside the simplified ETR band below 15 per cent. & Not structurally identified from the 2017-2022 public CbCR panel, which is pre-reform and does not observe a post-SESH density spike. & Post-SESH MNE-jurisdiction microdata with simplified ETRs, safe-harbour elections and enough support around the 15 per cent threshold to estimate missing and excess mass. \\
Channel split & Diagnostic only. The short-run split uses the Heckemeyer-Overesch meta-regression distinction between transfer-pricing or licensing and financial-planning channels. & The channel split does not affect top-up tax revenue in the legal allocation matrix. & Royalty, licence, service-fee, interest, intra-group debt and IP ownership data by MNE group and affiliate. \\
Substance relocation & Set to zero in the short-run behavioural response because payroll and tangible assets are held fixed in the SBIE calculation. & Real activity relocation is a longer-run margin and cannot be separated from book-profit shifting in aggregate CbCR. & Firm-level employment, payroll, tangible-asset and investment responses around anti-avoidance or GMT implementation shocks. \\
Destination matrix & Uses the Garcia-Bernardo and Jansky activity-weighted formula as a profit-conservation diagnostic. & Public CbCR does not observe origin-destination flows of shifted profit, so PPML or conditional logit would describe profit stocks rather than identify counterfactual relocation flows. & Origin-destination royalty, interest, service-fee or profit-flow data, plus destination choice sets, statutory and effective tax rates, withholding taxes, treaty links and anti-avoidance rules. \\
\bottomrule
\end{tabular}
}

\vspace{0.5em}
\begin{minipage}{0.96\textwidth}
\footnotesize\raggedright \textit{Notes:} Source: author's calculations based on \path{data/processed/diagnostics/step5_stage3_behavioural_identification_design.csv}.
\end{minipage}
\end{table}

